\begin{proof}
    In assumptions (1) through to (3), local properties placed on $\encoder( - )$ and $\decoder( - )$ in terms of generators and relations.
    These local properties lift to global properties as follows.
    \begin{enumerate}
    \item \textbf{Derivation (1)}.
          Assumption (1) states that $\sem{\encoder(x)}_\Gamma = \sem{x}_\Sigma$ for all $x \in \Sigma$.
          This means that $\sem{-}_\Gamma \circ \encoder \circ i = \sem{-}_\Sigma \circ i$.
          Then $\sem{-}_\Gamma \circ \encoder = \sem{-}_\Sigma$ by the universal property of free prop categories.
    \item \textbf{Derivation (2)}.
          Let $( X, Y ) \in Q$.
          By assumption (3), there exists an $X^* \in \decoder( X )$ and $Y^* \in \decoder( Y )$ such that $X^* \approx_E Y^*$.
          Since $X^* \in \decoder( X )$, then $\pi_E( X^* ) = \decoder( X )$.
          Since $Y^* \in \decoder( Y )$, then $\pi_E( Y^* ) = \decoder( Y )$.
          Since $X^* \approx_E Y^*$, then $\pi_E( X^* ) = \pi_E( Y^* )$.
          Then in conclusion, $\decoder( X ) = \pi_E( X^* ) = \pi_E( Y^* ) = \decoder( Y )$.
          Since $( X, Y )$ was an arbitrary relation in $Q$, then there exists a unique controlled prop functor $\decoder_Q: P_C( \Gamma ) / Q \to P( \Sigma ) / E$ such that $\decoder_Q \circ \pi_Q = \decoder$.
    \item \textbf{Derivation (3)}.
          Let $x \in \Sigma$.
          Then by assumption (2), there exists a $f \in \decoder( \encoder( g ) )$ such that $x \approx_E f$.
          Since $f \in \decoder( \encoder( x ) )$, then $\pi_E( f ) = \decoder( \encoder( x ) )$.
          Since $x \approx_E f$, then $\pi_E( x ) = \pi_E( f )$.
          Then $\pi_E( x ) = \pi_E( f ) = \decoder( \encoder( x ) )$.
          Since $x$ was as arbitrary generator, then $\pi_E \circ i = \decoder \circ \encoder \circ i$.
          Then $\pi_E = \decoder \circ \encoder$ by the universal property of prop categories.
    \end{enumerate}
     Next, it will be shown that $( \Sigma, E )$ is complete with respect to $\sem{-}_\Sigma$.
     Let $f, g \in P( \Sigma )$ such that $\sem{f}_\Sigma = \sem{g}_\Sigma$.
     Then $\sem{\encoder(f)}_\Gamma = \sem{f}_\Sigma = \sem{g}_\Sigma = \sem{\encoder(g)}_\Gamma$ by derivation (1).
     Since $( \Gamma, Q )$ is complete with respect to $\sem{-}_\Gamma$, then $\encoder( f ) \approx_Q \encoder( g )$.
     This means that $\pi_Q( \encoder( f ) ) = \pi_Q( \encoder( g ) )$, which implies that $\decoder_Q( \pi_Q( \encoder( f ) ) ) = \decoder_Q( \pi_Q( \encoder( g ) ) )$.
     Then $\decoder( \encoder( f ) ) = \decoder( \encoder( g ) )$ by derivation (2).
     Since $\pi_E = \decoder \circ \encoder$ by derivation (3), then $\pi_E( f ) = \pi_E( g )$.
     Then $f \approx_E g$.
     Since $f$ and $g$ were arbitrary, then $( \Sigma, E )$ is complete with respect to $\sem{-}_\Sigma$.
\end{proof}
